\documentclass[10pt,a4paper]{amsart}
\usepackage[margin=19mm]{geometry}
\usepackage{amsmath,amssymb,mathtools}
\usepackage[scr=rsfs,cal=euler]{mathalfa}
\usepackage{xfrac}
\usepackage[unicode,hidelinks,bookmarks=false]{hyperref}
\newcommand{\ie}{i.e.\ }
\newcommand{\cf}{cf.\ }
\newcommand{\resp}{respectively\ }
\newcommand{\Subsection}{\subsection}
\providecommand{\nobreakdash}{\nobreak\hbox{-}}
\setlength{\emergencystretch}{2em}
\multlinegap0pt
\usepackage{bm}
\usepackage{bbm}
\usepackage{dsfont}
\usepackage{xspace}
\usepackage{cancel}
\usepackage{mathtools}
\usepackage{amsthm}
\makeatletter
\@ifundefined{definition}{\newtheorem{definition}{Definition}}{}
\@ifundefined{lemma}{\newtheorem{lemma}[definition]{Lemma}}{}
\@ifundefined{theorem}{\newtheorem{theorem}[definition]{Theorem}}{}
\makeatother
\newcommand{\norm}{\mathrel{\unlhd}}
\title[A conjecture related to the nilpotency of groups with isomor]{A conjecture related to the nilpotency of groups with isomorphic non-commuting graphs}
\author{Valentina Grazian, Carmine Monetta}
\date{Source: arXiv:2302.01770v3 (CC BY 4.0)}
\begin{document}
\maketitle

\noindent\emph{Source and licence.} A conjecture related to the nilpotency of groups with isomorphic non-commuting graphs. Version arXiv:2302.01770v3 (CC BY 4.0), distributed under the Creative Commons Attribution 4.0 International licence, as recorded in the arXiv licence metadata for this exact version and shown on the abstract page https://arxiv.org/abs/2302.01770v3. Full rights notice: licenses/arxiv-2302.01770-CC-BY-4.0.txt.\par\smallskip

\noindent\textit{Editorial scope.} Counted unit: the Theorem whose source label is \texttt{exception} in the article (source0.tex lines 492--500, source label \texttt{exception}), delivered together with the article's own complete original proof of it (source0.tex lines 503--657, one \texttt{proof} environment of 155 lines). No mathematics is written, repaired or re-derived: statement and proof are the article's own text line for line. Required context retained uncounted: lem:p-groups (lines 334--344). Every definition, display and label of those lines is retained unchanged. Omitted from this package and not redistributed: 1--333, 345--491, 501--502, 658--670. Nothing inside a retained excerpt is omitted or altered: every retained body is delivered exactly as the article's lines are written, so no omission receipt of any kind accompanies this package. Citations: the retained excerpts carry no citation command at all, so no bibliography entry of the article is transcribed and no reference is redistributed. Reference adaptation: none was needed and none was made; every reference inside the delivered unit resolves inside it, so no unresolved reference or citation is produced. Printed slips kept literal: no printed word, symbol, hypothesis or number of the counted statement or of its proof was altered. Layout and class: the article's own class is \texttt{amsart} with packages including inputenc, babel, amsmath, amsthm, amssymb, amsfonts, mathrsfs, latexsym, textcomp, enumerate, ulem, xcolor; the wrapper is the lane's pinned \texttt{amsart} editorial wrapper at 10pt on A4, which restates the theorem-like environments the retained text needs and adds the article's own macro definitions that the retained text needs (\texttt{\textbackslash norm}), with extra wrapper packages (bm, bbm, dsfont, xspace, cancel, mathtools). The delivered numbering is the unit's own. Assets: no retained span contains a figure environment, a graphics-inclusion command, a PDF-page inclusion, a table or a TikZ picture; no image file is copied into this package, the package declares no asset dependency and no image file is redistributed with it. Licence: version arXiv:2302.01770v3 of this article is distributed under the Creative Commons Attribution 4.0 International licence, as recorded in the arXiv licence metadata for this exact version and shown on the abstract page https://arxiv.org/abs/2302.01770v3. A full rights notice is delivered at licenses/arxiv-2302.01770-CC-BY-4.0.txt. Characters: every retained excerpt is pure ASCII, so the delivered engine, XeLaTeX with its default Latin Modern text font, reports no missing character.
\begin{lemma}\label{lem:p-groups}
Let $P$ be a finite non-abelian AC-group and suppose that $P$ is a $p$-group for some prime $p$. Let $x \in Z_2(P) \backslash Z(P)$. Set $[P \colon Z(P)] = p^n$ and $[C_P(x) \colon Z(P)] = p^s$. If $c$ denotes the nilpotency class of $P$, then
\begin{enumerate}
\item $P' \leq C_P(x) \norm P$ (in particular $P'$ is abelian);
\item if $c > 2$ then $C_P(x)$ is the unique normal centralizer of $P$ and its order is maximum among the non-central element-centralizers of $G$ (in particular $Z_2(P) \leq C_P(x) = C_P(Z_2)$);
\item if $[P \colon C_P(x)] \geq p^2$ then $P/Z(P)$ has exponent $p$;
\item if $[P \colon C_P(x)] \geq p^2$ then $c \leq p$ and $Z_i(P) \leq C_P(x)$ for every $1 \leq i \leq c-1$;
\item if $y \in P \backslash C_P(x)$ and $[C_P(y) \colon Z(P)] = p^t$ then $n\geq s+t$ (in particular if $c > 2$ then $n \geq 2t$).
\end{enumerate}

\end{lemma}

\begin{theorem}\label{exception} 
Suppose $H$ is of type (3) with $[F \colon Z(H)] = q^f$, where $F/Z(H)$ is the kernel of $H/Z(H)$, $q$ is a prime and $f\geq 1$ is an integer. Then
\begin{enumerate}
\item $p \neq q$, that is, the Sylow $q$-subgroups of $G$ are abelian (possibly trivial);
\item $|Z(H)| > |Z(G)|$;
\item $[P \colon Z(P)] > p^4$;
\item none of the maximal subgroups of $G$ is abelian (in particular $P$ does not have maximal nilpotency class). 
\end{enumerate}
\end{theorem}

\begin{proof}
Set $|Z(H)| = bq^u$ for some integers $b, u \geq 1$ with $(b,q)=1$, so $|F| = bq^{u+f}$.  
Let $K/Z(H)$ denote a Frobenius complement of the group $H/Z(H)$. Then 
$|K| = [K \colon Z(H)]|Z(H)| = cbq^u$, for some integer $c\geq 2$ such that  $c$ divides $q^f -1$ (and so in particular $(c,q) = 1$). 
Take $k \in K \backslash Z(H)$, so $K = C_H(k)$, and $x_1, \dots, x_v \in F \backslash Z(H)$, one for each size of $H$-conjugacy classes, so $C_H(x_i) = C_F(x_i)$. In particular 
\[|C_H(k)| = |K| = c|Z(H)| \quad \text{ and } \quad |C_H(x_i)| = q^{t_i}|Z(H)|\] for distinct integers $t_i\geq 1$. Note that from $(c,q)=1$ we see that $|C_H(k)| \neq |C_H(x_i)|$.

Aiming for a contradiction, suppose there exists an isomorphism of graphs
$\phi \colon \Gamma_H \rightarrow \Gamma_G$. Let $g= \Phi(k)$, $M=C_G(g)$ and $y_i=\phi(x_i)$ for every $1 \leq i \leq v$. Note that $G/Z(G)$ is a $p$-group, so we can set $[M \colon Z(G)] = p^m$  for some integer $m\geq 1$ and $[C_G(y_i) \colon Z(G)]=p^{s_i}$ for some distinct integers $s_i\geq 1$.
Also, $|K| - |C_H(x_i)| = |M| - |C_G(y_i)|$, so $|M| \neq |C_G(y_i)|$, that is, $m \neq s_i$ for every $i\geq 1$.

\begin{description}
\item[Claim 1] $p \neq q$, so part (1) holds.

\smallskip
\noindent
Proof: 
Since $\Gamma_H \cong \Gamma_G$ we have
\begin{equation}\label{eq2}|H| - |K| = |G| - |M| \Longrightarrow
%|K|([H \colon K] - 1) = |M|([G \colon M] - 1) \Longrightarrow
q^ubc(q^f- 1) = p^{m+r}a(p^{n-m}- 1)\end{equation}

\vspace{0.2cm}
and
\begin{equation}\label{eq3} |C_H(x_i)| - |Z(H)| = |C_G(y_i)| - |Z(G)| \Longrightarrow
%|Z(H)|([C_H(x_i) \colon Z(H)] - 1) = |Z(G)|([C_G(y_i) \colon Z(G)] - 1) \Longrightarrow
q^ub(q^{t_i} - 1) = p^{r}a(p^{s_i}- 1) .\end{equation}

Aiming for a contradiction, suppose $p=q$. Then $(p,b)=(p,c)=(p,a)=1$. From Equation (\ref{eq2}) we deduce that $p^ u = p^{m+r}$ and  from Equation (\ref{eq3}) we get $p^ u = p^{r}$, a contradiction. Therefore $p\neq q$.

\bigskip
\item[Claim 2] for every $d\leq r$ we have that $p^d$ divides $bc$ if and only if $p^d$ divides $b$.

\smallskip
\noindent
Proof: 
By assumption $p^d$ divides $p^r$ and by Claim $1$ we have $p\neq q$. We conclude by the following equality:
\begin{equation}\label{eq4} |K| - |Z(H)| = |M| - |Z(G)| 
\Longrightarrow  q^u( bc - b ) = p^ra(p^m - 1)\end{equation}



\bigskip
\item[Claim 3] $p^{r+1}$ does not divide $bc$ and $q^{u+1}$ does not divide $a$. In particular $(p,c)=1$ and so $(p,[H \colon Z(H)])=1$.

\smallskip
\noindent
Proof:  Recall that $|K| - |C_H(x_i)| \neq 0 $ since $c$ does not divide $q$,  and so  $|M| \neq |C_g(y)|$. Note that
\begin{equation}\label{eq5} |K| - |C_H(x_i)| = |M| - |C_G(y_i)| \Longrightarrow 
q^u(bc - bq^{t_i}) = p^ra(p^m - p^{s_i}). \end{equation}


From Claim 1 we know $p\neq q$ and by Equation (\ref{eq5}) we get that $p^{r+1}$ divides $(bc - bq^{t_i})$. 
Aiming for a contradiction, suppose $p^{r+1}$ divides $bc$. Then $p^{r+1}$ must divide $b$ and by Equation (\ref{eq4}) we deduce that $p^{r+1}$ divides $p^ra(p^m-1)$, a contradiction. Therefore $p^{r+1}$ does not divide $bc$.

Similarly, if $q^{u+1}$ divides $a$,  then Equation (\ref{eq5}) implies that $q$ divides $bc$, reaching again a contradiction.

Finally, if $p$ divides $c$, then $p$ does not divide $(c-1)$ and by Equation (\ref{eq4}) we deduce that $p^r$ divides $b$. But then $p^{r+1}$ divides $bc$, a contradiction.


\bigskip
\item[Claim 4] We have
\begin{equation}\label{eqC} p^{m}(p^{n-m}- 1)(c - 1)= (p^{m}- 1)c(q^f - 1).\end{equation} In particular 
\begin{enumerate}
\item $n > 2m$;
\item $c < p^m < q^f < p^{n-m}$; and
\item $|Z(H)| > |Z(G)|$, so part (2) holds.
\end{enumerate}

\smallskip
 \noindent
Proof: From Equation (\ref{eq2}) we get
\[ |Z(H)| = q^ub = \frac{p^{r+m}a(p^{n-m}- 1)}{c(q^f - 1)}.\]
Substituting this into Equation (\ref{eq4}) we obtain
%\[ \frac{p^{r+m}a(p^{n-m}- 1)}{c(q^f - 1)}\cdot (c - 1)= p^{r}a(p^{m}- 1) \]
%and so
$$ p^{m}(p^{n-m}- 1)(c - 1)= (p^{m}- 1)c(q^f - 1),$$
as wanted.

Recall that $c$ divides $q^f-1$, so $c \leq q^f-1$. Aiming for a contradiction, suppose $c=q^f-1$. Then Equation (\ref{eqC}) becomes
\[ p^{m}(p^{n-m}- 1)(c - 1)= (p^{m}- 1)c^2 , \]
that we can rewrite as
\[(p^n- p^m(1+c) +c)c = p^{m}(p^{n-m}- 1). \]

In particular $p^m$ divides $(p^n- p^m(1+c) +c)c $. However, $p^m$ divides $p^n$ while $(p,c)=1$ by Claim 3, and we reach a contradiction. Thus $c<q^f-1$.

As a consequence, we have $c-1 \leq q^f-1$ and from Equation (\ref{eqC}) we deduce that we must have $p^{m}(p^{n-m}- 1) \geq (p^{m}- 1)c$ and so 
\[c \leq \frac{p^{m}(p^{n-m}- 1)}{ (p^{m}- 1)}. \]

Rewrite Equation (\ref{eqC}) as
$$(p^{n-m}- 1)(p^mc - p^m)= (p^{m}c- c)(q^f - 1)$$
and notice that
$$c > p^m \Leftrightarrow p^mc - p^m > p^mc - c \Leftrightarrow p^{n-m}-1 < q^f-1
\Leftrightarrow p^{n-m} < q^f.$$  

Aiming for a contradiction suppose $n \leq 2m$. Then $p^{n-m} \leq p^m$ and by Equation (\ref{eqC}) $p^m$ divides $q^f-1$, so $p^{n-m} \leq p^m \leq q^f-1$.
Hence by what we noted above we obtain $c \geq p^m$. On the other hand, $c \leq \frac{p^{m}(p^{n-m}- 1)}{ (p^{m}- 1)} \leq p^m$. Thus the only option is $c=p^m$, contradicting Claim 3. Therefore $n > 2m$.

Now, if $c > p^m$ then $q^f > p^{n-m} > p^m$ contradicting the fact that by Equation (\ref{eqC}) $p^m$ divides $q^f-1$. Therefore $c < p^m < q^f < p^{n-m}$. Finally, from Equation (\ref{eq4}) we conclude that $|Z(H)| > |Z(G)|$.


\bigskip
\item[Claim 5] for every $i\geq 1$ we have $n > 2s_i$.

\smallskip
\noindent
Proof: 
Note that, using once again the fact that $\Gamma_H \cong \Gamma_G$, we have

\begin{equation}\label{eq6} |H| - |C_H(x_i)| = |G| - |C_G(y_i)| \Longrightarrow 
q^{u+t_i}b(q^{f-t_i}c - 1) = p^{r+s_i}a(p^{n-s_i} - 1). \end{equation}

From Equation (\ref{eq6}) we get
\[ |Z(H)| = q^ub = \frac{p^{r+s_i}a(p^{n-s_i}- 1)}{q^{t_i}(q^{f-t_i}c - 1)}.\]
Substituting this into Equation (\ref{eq1}) we obtain
\[ \frac{p^{r+s_i}a(p^{n-s_i}- 1)}{q^{t_i}(q^{f-t_i}c - 1)}\cdot (q^fc - 1)= p^{r}a(p^{n}- 1) \]
and so 
\[ p^{s_i}(p^{n-s_i}- 1)(q^fc - 1) = (p^{n}- 1)q^{t_i}(q^{f-t_i}c - 1). \]
By Claim 1, $p\neq q$. So we deduce that $p^{s_i}$ divides $q^{f-t_i}c - 1$. In particular, 
\begin{equation}\label{last}
p^{s_i} \leq q^{f-t_i}c - 1 \leq q^{f-t_i}c.
\end{equation}

Since $\Gamma_H \cong \Gamma_G$, we also obtain
\begin{equation}\label{eq1} |H| - |Z(H)| = |G| - |Z(G)| \Longrightarrow 
q^{u}b(q^fc - 1) = p^{r}a(p^{n} - 1). \end{equation}
From Equation (\ref{eq3}) we get
\[ |Z(H)| = q^ub = \frac{p^{r}a(p^{s_i}- 1)}{q^{t_i} - 1}.\]
Substituting this into Equation (\ref{eq1}) we obtain
%\[ \frac{p^{r}a(p^{s_i}- 1)}{q^{t_i} - 1}(q^fc- 1) = p^{r}a(p^{n}- 1) \]
%and so 
\[ q^fc = \frac{(p^{n}- 1)(q^{t_i} - 1)}{p^{s_i}- 1} + 1. \]

Now, $p^n - 1 = (p^s_i - 1)p^{n-s_i} +(p^{n-s_i} - 1)$ and so we deduce  

\[\begin{aligned} q^fc &= \frac{((p^s_i - 1)p^{n-s_i} +(p^{n-s_i} - 1))(q^{t_i} - 1)}{p^{s_i}- 1} + 1 
\\ &= p^{n-s_i}q^{t_i}  + \frac{(p^{n-s_i} - 1)(q^{t_i} - 1)}{p^{s_i}- 1} + 1 - p^{n-s_i}
\\ &= p^{n-s_i}q^{t_i} - \frac{(p^{s_i}-q^{t_i})(p^{n-s_i}-1)}{p^{s_i} -1}.
\end{aligned} \]


%\\ &= p^{n-s_i}q^{t_i}  + \frac{p^{n-s_i}q^{t_i} - p^{n-s_i} -q^{t_i} +1 + p^{s_i}- 1 - p^n +p^{n-s_i}}{p^{s_i}- 1}

By Claim 4 we have $|Z(H)| > |Z(G)|$, so by Equation (\ref{eq3}) we conclude that $p^{s_i} > q^{t_i}$. Thus $\frac{(p^{s_i}-q^{t_i})(p^{n-s_i}-1)}{p^{s_i} -1} > 0$ and $q^fc < p^{n-s_i}q^{t_i}$. Combining this with (\ref{last}) we get
\[ p^{s_i}q^{t_i} \leq  q^fc < p^{n-s_i}q^{t_i}.\]
Therefore $p^{s_i} < p^{n-s_i}$, implying $n > 2s_i$.
\end{description}

Now, to prove parts (3) and (4) of the statement, recall that for every $i\geq 1$ we have $m \neq s_i$ and Claims 4 and 5 give $n > 2m$ and $n > 2s_i$. Thus we conclude that $n >4$ and $s_i \neq n-1 \neq m$, implying that the maximal subgroups of $G$ are not abelian. In particular, by Lemma \ref{lem:p-groups}(2) we conclude that if $P$ has nilpotency class greater than $2$ then $C_P(Z_2(P))$ is not a maximal subgroup of $P$. Hence $P$ does not have maximal nilpotency class.

%\item[Claim 8] \textcolor{red}{If $p=2$ then for every $y \in G \backslash Z(G)$ we have $[C_G(y) \colon Z(G)]\geq p^2=4$. In particular, $[G \colon Z(G)]\geq 2^7$.}

%Proof: Note that the statement is equivalent to prove that $m\neq 1$ and $s_i \neq 1$ for every $y_i=\phi(x_i)$ with $x_i \in F\backslash Z(H)$. The fact that $m>1$ follows from Claim 4 (2), since $1\neq c < 2^m.$ Now, aiming for a contradiction, suppose $s_i = 1$ for some $y_i=\phi(x_i)$ with $x_i \in F\backslash Z(H)$. Since $|Z(H)|> |Z(G)|$ by Claim 4 (3), Equation (\ref{eq3}) implies $q^{t_i} < p^{s_i} = 2$, a contradiction. Therefore we must have $m,s_i \geq 2$ for every $i$. Note that $m\neq s_i$ by Equation (\ref{eq6}), so at least one among $m$ and $s_i$ is greater than $2$.  By Claim 4(1) and Claim 7 we deduce that $n > 2 \cdot \max{m,s_i} \geq 6$, hence $n \geq 7$, as wanted. 

\end{proof}

\end{document}
